\documentclass[11pt,letterpaper]{article}
 
% ---------- Packages ----------
\usepackage[margin=1in]{geometry}
\usepackage[T1]{fontenc}
\usepackage{lmodern}
\usepackage{microtype}
\usepackage{amsmath,amssymb}
\usepackage{booktabs}
\usepackage{graphicx}
\usepackage{float}
\usepackage{caption}
\usepackage{enumitem}
\usepackage{xcolor}
\usepackage{titlesec}
\usepackage{fancyhdr}
\usepackage[hidelinks]{hyperref}
\usepackage{siunitx}
 
\definecolor{gatorblue}{HTML}{0021A5}
\definecolor{gatororange}{HTML}{FA4616}
 
\titleformat{\section}{\large\bfseries\color{gatorblue}}{\thesection}{0.6em}{}[\color{gatororange}\titlerule]
\titleformat{\subsection}{\normalsize\bfseries}{\thesubsection}{0.6em}{}
\setlist{itemsep=2pt, topsep=4pt}
 
\pagestyle{fancy}
\fancyhf{}
\cfoot{\small\thepage}
\renewcommand{\headrulewidth}{0pt}
 
% ---------- Metadata ----------
\newcommand{\NoteTitle}{Algorithmic Genetic Feature Selection to Improve XGBoost on E-mini S\&P~500 Order-Book Data.}
\newcommand{\Authors}{Manny Osorio, Emily Diaz-Silva, Benjamin Friedman}
\newcommand{\NoteDate}{\today}
\newcommand{\TeamName}{}
 
\begin{document}
 
% ---------- Title block ----------
\begin{center}
    {\LARGE\bfseries\color{gatorblue} \NoteTitle}\\[6pt]
    {\large Quant Note -- Systematic Trading Track}\\[4pt]
    \Authors
\end{center}
 
% ============================================================
\section{Summary}
% ============================================================
\textbf{Overall Strategy:} On each ES order-book update, we combine GA-selected book features and XGBoost outputs into a standardized score. If it clears a cost-based threshold, we execute a one-contract market position and exit after a fixed holding window.
 
We evaluate whether the ten-level order book of the CME E-mini S\&P~500 future (ES) predicts mid-price direction over the next [$h$] seconds better than a top-of-book rule. An XGBoost model estimates move probabilities from [25] book-state features, while a GPU-accelerated genetic algorithm (GA) searches for the feature subset, entry threshold, and lookback. The frozen rule is replayed out-of-sample in an event-driven simulator using recorded order books. Net of fees, the strategy achieves [X] trades with a mean net P\&L of \$[Y] per trade, a net Sharpe ratio of [S], and a maximum drawdown of [D]\% over [N] days, comparing against random and level-1 baselines.
 
% ============================================================
\section{Economic and Data Hypothesis}
% ============================================================
 
\subsection{Economic Hypothesis}
The ten-level ES order book contains microstructural information regarding mid-price direction over $h$ seconds. Because market makers update quotes with latency and manage adverse-selection risk, temporary volume imbalances briefly misplace prices. Aggressive traders and stale market makers absorb these costs. Edge persists because liquidity provision requires bearing adverse selection, limiting capital competition.
 
\paragraph{Testable prediction.} Net of spread and fees, the frozen rule earns positive out-of-sample returns; realized returns scale across buckets of the model's directional score; and the rule beats random feature subsets and level-1 imbalance rules.
 
\paragraph{It fails if.} (1) Out-of-sample net returns are negative after transaction costs; (2) realized returns are non-monotone relative to directional scores; or (3) performance fails to beat baselines.
 
% ============================================================
\section{Data and Universe}
% ============================================================
We trade a single instrument: the CME E-mini S\&P~500 front-month future (\texttt{ES.c.0}, ESH4 contract). Market data comes from Databento~\cite{Databento} (\texttt{GLBX.MDP3}, \texttt{mbp-10} schema) at nanosecond resolution. The sample spans [2024-01-01 to 2024-01-31], containing [N] book events. 

We utilize \texttt{ts\_recv} timestamps. Crossed/locked books are dropped, empty levels are padded with zeros, and maintenance breaks (17:00--18:00 ET) are excluded.
 
% ============================================================
\section{Methodology}
% ============================================================
\subsection{Model Training and Genetic Search}
\paragraph{XGBoost.} The gradient-boosted tree classifier~\cite{chen2016xgboost}
(GPU histogram method, \texttt{multi:softprob}, three classes) outputs
$[P_{\text{sell}},P_{\text{hold}},P_{\text{buy}}]$. Out-of-fold probabilities are produced with an expanding-window purged walk-forward scheme. The directional score is $q_i = P_{\text{buy}}(\mathbf{x}_i) - P_{\text{sell}}(\mathbf{x}_i) \in [-1,1]$.
 
\paragraph{Genome.} Encoded as a 64-bit chromosome (Table~\ref{tab:genome}) governing feature flags, Gray-coded lookback, entry threshold, and risk parameters.
 
\begin{table}[h]
\centering
\caption{Chromosome layout (64 bits).}
\label{tab:genome}
\small
\begin{tabular}{@{}lcl@{}}
\toprule
\textbf{Field} & \textbf{Bits} & \textbf{Meaning} \\
\midrule
Feature flags & 28 & Activates GA input column $j$ \\
Lookback & 12 & Initial events skipped \\
Threshold & 12 & Standardized entry cutoff ($\tau$) \\
Risk & 12 & Position scaling parameter \\
\midrule
\textbf{Total} & \textbf{64} & \\
\bottomrule
\end{tabular}
\end{table}
 
\paragraph{Rule and fitness.} Standardized signals form $S_i = \sum_{j\in\mathcal{S}_b} z_{i,j}$, yielding positions $p_i \in \{+1, -1, 0\}$. A fused Triton kernel~\cite{tillet2019triton} evaluates populations in GPU memory, maximizing a skew-penalized Sharpe score $\mathcal{F}(b)$.
 
\paragraph{Evolution.} Populations evolve across GPU islands using elitism, crossover, and mutation, with periodic elite migration.
 
\paragraph{Adaptive Re-Evolution and Validation} Re-evolution occurs across rolling windows under purged walk-forward constraints to handle non-stationarity without data leakage.

\subsection{Strategy}
\paragraph{Signal construction.} Evaluated at every depth update using frozen in-sample weights, feature masks, and thresholds after a warm-up period.
 
\paragraph{Portfolio Construction}
Single-instrument execution (ESH4). Long, short, or flat based on $S_t$ thresholds from flat positions. Exposure does not stack, and positions flatten before maintenance breaks.
 
\paragraph{Position Sizing}
Fixed size of $n=1$ contract per trade to isolate model alpha.
 
\paragraph{Rebalancing Frequency} Event-driven updates on book snapshots, throttled by a 500\,ms minimum interval and exit cooldown. Positions close after holding period $h$ or before maintenance breaks.
 
\paragraph{Execution Assumptions}
Market orders fill against simulated L2 depth. A fixed latency of $[\,]$\,ms is applied, alongside a \$2.00 per-contract fee.
 
\subsection{Backtest Design}
Evaluated via NautilusTrader~\cite{nautilus} using historical MBP-10 data. All tuning occurs strictly out-of-engine on in-sample splits.
 
% ============================================================
\section{Results}
% ============================================================
\subsection{Performance Summary}
\begin{table}[H]
\centering
\caption{Performance statistics (net of costs)}
\begin{tabular}{lrrr}
\toprule
Metric & In-Sample & Out-of-Sample & Benchmark \\
\midrule
Annualized return (\%)    & \dots & \dots & \dots \\
Annualized volatility (\%) & \dots & \dots & \dots \\
Sharpe ratio              & \dots & \dots & \dots \\
Sortino ratio             & \dots & \dots & \dots \\
Max drawdown (\%)         & \dots & \dots & \dots \\
Calmar ratio              & \dots & \dots & \dots \\
Hit rate (\%)             & \dots & \dots & \dots \\
Turnover (\% / month)     & \dots & \dots & \dots \\
\bottomrule
\end{tabular}
\end{table}
 
\subsection{Cumulative Performance}
\begin{figure}[H]
\centering
\fbox{\parbox{0.8\linewidth}{\centering\vspace{3cm}Equity curve placeholder\vspace{3cm}}}
\caption{Cumulative net returns vs.\ benchmark.}
\end{figure}
 
\subsection{Robustness and Attribution}
\dots
 
% ============================================================
\section{Risk Management}
% ============================================================
\subsubsection{Risk Controls}
A rule-based risk manager overrides orders without changing size:
\begin{itemize}
\item \textbf{Pre-trade gates:} Blocks trades during maintenance breaks, major macro releases (FOMC, CPI), roll weeks, wide spreads, or extreme volatility regimes.
\item \textbf{In-trade exits:} Enforces price stops, time stops ($h$), and pre-break flat triggers.
\item \textbf{Session limits:} Halts trading on daily losses exceeding 2\%, equity peak drawdowns of 5%, or sequential loss streaks.
\end{itemize}
 
% ============================================================
\section{Liquidity and Capacity}
% ============================================================
\subsection{Liquidity Profile}
\dots
 
\subsection{Transaction Costs and Market Impact}
\dots
 
\subsection{Capacity Estimate}
\begin{table}[H]
\centering
\caption{Net Sharpe ratio vs.\ assumed AUM}
\begin{tabular}{lrrrr}
\toprule
AUM (\$M) & 1 & 10 & 50 & 100 \\
\midrule
Net Sharpe & \dots & \dots & \dots & \dots \\
Max ADV participation (\%) & \dots & \dots & \dots & \dots \\
\bottomrule
\end{tabular}
\end{table}
 
% ============================================================
\section{Limitations and Next Steps}
% ============================================================
\subsection{Limitations}
\begin{itemize}\setlength{\itemsep}{0pt}\setlength{\parskip}{0pt}
    \item \textbf{Data:} trained on one month of ES only, which captures a single market regime; MBP-10 has no queue position.
    \item \textbf{Search budget:} compute limited the GA to 100 generations, so fitness may not have converged and results may be a local optimum.
    \item \textbf{Overfitting:} many parameter sets tested on overlapping returns from a short sample, so results may not generalize.
    \item \textbf{Regime risk:} untested under stress; bursts may not revert in high-volatility markets.
    \item \textbf{Execution:} assumes 1.32-tick round trips with no impact or latency model.
\end{itemize}

\subsection{Next Steps}
\begin{itemize}\setlength{\itemsep}{0pt}\setlength{\parskip}{0pt}
    \item \textbf{More data and generations:} ideally train on 2+ years, including 2022 and Aug.\ 2024, and run the GA until fitness plateaus; test on NQ.
    \item \textbf{Validation:} hold out a clean test period and fix overlapping fitness.
    \item \textbf{Execution:} simulate passive fills with MBO data, then paper-trade live in MES.
\end{itemize}
 
\addcontentsline{toc}{section}{References}
\begin{thebibliography}{99}
\small

\bibitem{chen2016xgboost}
T.~Chen and C.~Guestrin,
``XGBoost: A Scalable Tree Boosting System,''
in \textit{Proceedings of the 22nd ACM SIGKDD International Conference on Knowledge Discovery and Data Mining}, 2016, pp. 785--794.
\url{https://doi.org}

\bibitem{christodoulou2025}
C.~N. Christodoulou-Volos,
``Optimizing Trading Strategies Using Genetic Algorithms: A Review and Implementation,''
\textit{Edelweiss Applied Science and Technology}, vol.~9, no.~10, pp.~334--345, 2025.
\url{https://doi.org/10.55214/2576-8484.v9i10.10407}

\bibitem{bailey2014dsr}
D.~H. Bailey and M.~L\'opez de Prado,
``The Deflated Sharpe Ratio: Correcting for Selection Bias, Backtest Overfitting and Non-Normality,''
\textit{Journal of Portfolio Management}, vol.~40, no.~5, pp.~94--107, 2014.
\url{https://doi.org/10.3905/jpm.2014.40.5.094}

\bibitem{Databento}
Databento. (2026). GLBX.MDP3: CME Globex market data, MBP-10 schema. Databento.
\url{https://databento.com/catalog/cme/GLBX.MDP3}

\end{thebibliography}
 
\end{document}